\section{Proportion Logic and Relational Operators}
\label{sec:proportion_logic}

Proportion forms the bedrock of invariant structure. Consider two elements $A, B \in \mathcal{V}$ under a relational ratio operator.

\begin{definition}[Proportional Invariance]
For $A, B, C, D \in \mathcal{V}$ with non-zero denominators, equality of proportion is defined as:
\begin{equation}
A : B = C : D \iff \frac{A}{B} = \frac{C}{D}
\end{equation}
Under a scale transformation $\mathcal{S}_\lambda(x) = \lambda x$ for $\lambda \neq 0$, the ratio remains invariant:
\begin{equation}
\frac{\lambda A}{\lambda B} = \frac{A}{B}
\end{equation}
\end{definition}

This simple property demonstrates how structural information (the ratio $A/B$) is completely preserved even when absolute magnitudes ($A, B \to \lambda A, \lambda B$) undergo global transformation.

\subsection{The Reference / Balancing Third Term}
In formal logic and relational dynamics, two opposing or relative terms $(A, B)$ require a reference condition $C$ to form a balanced relational tuple:
\begin{equation}
(A, B; C)
\end{equation}
where $C$ acts as an invariant boundary condition, reference frame, or balancing mediator.
